%! AP Statistics — Unit 8, Lesson 8.5 — Group Activity
\documentclass[10pt]{article}
\usepackage{apstats-article}
\usepackage{apstats-boxes}
\newcommand{\TallMath}[1]{$\displaystyle #1\rule[-1.4em]{0pt}{3.2em}$}

\begin{document}

\pageheader{Unit 8, Lesson 8.5}{Group Activity}
\namepartnerperiod

\vspace{0.08in}

\begin{scenariobox}[Scenario A: Extracurricular Activities at Three Schools]{sky}
A researcher randomly selects 100 students from each of three high schools (Schools A, B, and C)
and asks which extracurricular activity (Sports, Arts, or Neither) each student participates in.
The results are shown below.

\vspace{0.08in}
\renewcommand{\arraystretch}{1.4}
\begin{tabularx}{0.82\linewidth}{l cccc}
    & \textbf{Sports} & \textbf{Arts} & \textbf{Neither} & \textbf{Row Total} \\
    \hline
    \textbf{School A} & 45 & 32 & 23 & 100 \\
    \textbf{School B} & 38 & 40 & 22 & 100 \\
    \textbf{School C} & 52 & 28 & 20 & 100 \\
    \hline
    \textbf{Col Total} & 135 & 100 & 65 & 300 \\
\end{tabularx}
\end{scenariobox}

\begin{enumerate}[label=\textbf{A\arabic*.}]

    \item \textbf{Identify the test type.} Is this a test for homogeneity or independence?
          Circle your answer and explain why based on the study design.
          \writelines{3}

    \item \textbf{Write the hypotheses.} State $H_0$ and $H_a$ in the context of this study.

          $H_0$: \writeline
          $H_a$: \writeline

    \item \textbf{Independence condition.} How was independence of observations achieved
          in this study? Was the 10\% condition necessary? Explain.
          \writelines{2}

    \item \textbf{Large-counts condition.} Calculate the expected counts for two cells:
          School A/Sports and School C/Neither. Then check whether the large-counts
          condition is plausibly met.

          Expected count (School A / Sports): \blank{5cm}

          Expected count (School C / Neither): \blank{5cm}

          Large-counts condition satisfied? \blank{2cm} \quad Explain: \writeline

\end{enumerate}

\vspace{0.1in}

\begin{scenariobox}[Scenario B: Political Affiliation and Policy Views]{sky}
A random sample of 300 adults is asked their political affiliation (Democrat, Republican,
or Independent) and their view on a new policy (Support or Oppose).
The results are shown below.

\vspace{0.08in}
\renewcommand{\arraystretch}{1.4}
\begin{tabularx}{0.72\linewidth}{l ccc}
    & \textbf{Support} & \textbf{Oppose} & \textbf{Row Total} \\
    \hline
    \textbf{Democrat}    & 85 & 15  & 100 \\
    \textbf{Republican}  & 30 & 70  & 100 \\
    \textbf{Independent} & 55 & 45  & 100 \\
    \hline
    \textbf{Col Total}   & 170 & 130 & 300 \\
\end{tabularx}
\end{scenariobox}

\begin{enumerate}[label=\textbf{B\arabic*.}]

    \item \textbf{Identify the test type.} Is this a test for homogeneity or independence?
          Circle your answer and explain why based on the study design.
          \writelines{3}

    \item \textbf{Write the hypotheses.} State $H_0$ and $H_a$ in the context of this study.

          $H_0$: \writeline
          $H_a$: \writeline

    \item \textbf{Independence condition.} How was independence of observations achieved
          in this study? Was the 10\% condition necessary? Explain.
          \writelines{2}

    \item \textbf{Large-counts condition.} Calculate the expected counts for two cells:
          Democrat/Support and Republican/Oppose. Then check whether the large-counts
          condition is plausibly met.

          Expected count (Democrat / Support): \blank{5cm}

          Expected count (Republican / Oppose): \blank{5cm}

          Large-counts condition satisfied? \blank{2cm} \quad Explain: \writeline

\end{enumerate}

\vspace{0.08in}

\begin{extensionbox}
\textbf{Tier E --- Extension:} Both Scenario A and Scenario B involve two categorical
variables appearing in a two-way table. A classmate says both scenarios call for a test
for independence. Is the classmate correct? Explain in 2--3 sentences what feature of
the study design determines which test to use.
\writelines{4}
\end{extensionbox}

\end{document}
